
Federated Learning (FL) is a distributed training paradigm in which multiple clients collaborate to build a global model without directly sharing their data, sending only model updates to the server. However, in adversarial settings, malicious or faulty clients can manipulate these updates (e.g., through Byzantine behavior and poisoning), compromising training convergence and degrading the performance of the final model. In this context, the main challenge is to define aggregation mechanisms that reduce the influence of anomalous updates while maintaining competitive performance under realistic conditions. This work aims to evaluate robust aggregation algorithms applied to FL under different levels of adversariality. To this end, robust aggregation methods—such as Krum and variants based on robust statistics (e.g., coordinate-wise median and trimmed mean)—will be compared against a reference mean-based aggregator. Performance will be analyzed using learning metrics (such as accuracy and F1-score), training stability, and sensitivity to the proportion of malicious participants, seeking to highlight the advantages, limitations, and applicability conditions of each aggregator. It is expected that the results will support the selection of robust aggregation strategies for FL systems exposed to untrusted participants.

\keywords{federated learning; robust aggregation; Byzantine failures; model poisoning; machine learning security}



%Alzheimer's Disease is an incurable neurodegenerative condition that presents itself as one of the main global public health challenges, affecting not only the patient but also the people around them. In this context, the main difficulty regarding the disease is its early detection, since several of its symptoms can be confused with normal aging. This work aims to validate speech-based audio signals as a promising biomarker to perform this detection. To that end, two classification models will be trained and validated—a Decision Tree and XGBoost—with a focus on using the SHAP technique to promote not only the accuracy but also the explainability of the models. The performance will be compared based on metrics such as accuracy, sensitivity, specificity, F1-Score, and the confusion matrix, as well as on the ability to explain which attributes collaborated to the final decision. It is expected that the obtained results will highlight the advantages and limitations of each approach, offering insights into the applicability of explainable models in the clinical context and thus providing better support for the diagnosis of Alzheimer's Disease.

% Separe as Keywords por ponto e vírgula.
%\keywords{Alzheimer's Disease; Speech Analysis; Machine Learning; Explainable AI; Decision Tree; XGBoost.}